\section{Data and methods}\label{sec:methods}

\subsection{Cities, tracts and population groups}

The study covers the \nCities{} municipalities with a 2022 Census population of at least 300,000, a set that
includes all 27 state capitals. The spatial units are 2022 census tracts (\emph{setores censit\'arios}). We keep
tracts classified as urban (\texttt{SITUACAO} = \emph{Urbana}) with at least one resident, \nTractsUrban{} tracts
holding \sharePopUrban\% of the population of the \nCities{} municipalities. We restrict the analysis to urban
tracts because several municipalities (e.g.\ Porto Velho or Manaus) contain very large rural tracts whose
centroids lie kilometres from any street, which would inflate network distances for reasons unrelated to urban
design. Population counts come from the Census tract aggregates for colour or race: \emph{branca}, \emph{preta},
\emph{amarela}, \emph{parda} and \emph{ind\'igena}. Cells suppressed for confidentiality are set to zero.

\subsection{Street networks}

For each municipality we download the OpenStreetMap pedestrian network (all ways usable on foot, including
footpaths and stairways) with OSMnx \citep{boeing2017osmnx} for the union of its urban tracts buffered by 2~km, so
that walking paths may leave the urban footprint. Graphs are simplified, projected to the SIRGAS~2000 UTM zone of
the city, made undirected, and reduced to their largest connected component. Each tract is represented by a
point guaranteed to lie inside the polygon and snapped to the nearest network node; the snap distance is kept and
added to network distances.

\subsection{Local environments and the segregation index}

Let $d_{ij}$ be the distance between tracts $i$ and $j$. We use two distances: Euclidean $d^{E}_{ij}$, and network
$d^{N}_{ij} = s_i + \ell(n_i, n_j) + s_j$, where $n_i$ is the node nearest to tract $i$, $s_i$ the snap distance and
$\ell$ the shortest-path length. Because the snapped path is itself a path between the two points,
$d^{N}_{ij} \geq d^{E}_{ij}$. Both distances enter the same triangular kernel,
\begin{equation}
  w_{ij} = \begin{cases} 1 - d_{ij}/b & \text{if } d_{ij} \le b,\\ 0 & \text{otherwise,}\end{cases}
  \qquad w_{ii} = 1,
  \label{eq:kernel}
\end{equation}
with bandwidth $b = 2{,}000$~m, as in \citet{knaap2024segregated}. The local environment of tract $i$ is the
spatially weighted population $\tilde{\tau}_{im} = \sum_j w_{ij}\tau_{jm}$ of each group $m$, and
$\tilde{\pi}_{im} = \tilde{\tau}_{im} / \sum_m \tilde{\tau}_{im}$. The entropy of each local environment is
$\tilde{E}_i = -\sum_m \tilde{\pi}_{im}\ln\tilde{\pi}_{im}$, and the spatial information theory index is
\begin{equation}
  H = 1 - \frac{1}{T E}\sum_i \tilde{t}_i \tilde{E}_i ,
  \label{eq:H}
\end{equation}
where $\tilde t_i$ is the local population, $T$ the total population and $E$ the city-wide entropy
\citep{reardon2004measures}. $H$ equals 0 when every local environment has the city's composition and 1 when every
local environment contains a single group. We compute $H$ over the five colour-or-race groups (a group absent from
a city contributes zero). As secondary measures we compute the spatial versions of dissimilarity, isolation and
entropy for the \emph{preta}+\emph{parda} population. All indices are computed with the PySAL \texttt{segregation}
package \citep{cortes2026segregation,cortes2020open}. We build the weighted local environments with a sparse
matrix product identical to the one the package applies to a \texttt{libpysal} weights object, which a unit test
verifies to machine precision, and then call the package's estimators. For each city we record
$H^{E}$, $H^{N}$, the gap $\Delta H = H^{N} - H^{E}$ and the relative gap $\%\Delta H = 100\,\Delta H / H^{E}$.

\paragraph{Population-matched environments.} Because $d^{N}_{ij} \ge d^{E}_{ij}$, a network environment of radius
$b$ always contains fewer (weighted) people than a Euclidean one, and smaller environments are, other things equal,
less diverse. Following the concern raised by \citet{saporito2026reconsidering}, we also compare $H^{N}$ with a
Euclidean index computed from \emph{population-matched} environments: for each tract we find, by bisection, the
Euclidean radius $b_i \le 2{,}000$~m at which the kernel-weighted population of its Euclidean environment equals
that of its 2-km network environment, and recompute $H$ with these tract-specific radii. Any remaining gap cannot
be attributed to the number of people in local environments.

\subsection{Network topology and configurational fragmentation}

Topology is measured on all network nodes lying inside the urban footprint (the union of populated urban
tracts). We keep every piece rather than only the largest component, because planned cities such as Palmas
consist of superblocks that connect through streets outside the footprint. Table~\ref{tab:topodef} lists the
measures, which follow \citet{knaap2024segregated}. Densities are expressed per km$^2$ of footprint.

\begin{table}[htbp]
\centering\small
\caption{Street-network topology measures ($v$ nodes, $e$ edges, $p$ connected components).}\label{tab:topodef}
\begin{tabular}{p{4.2cm}p{9.3cm}}
\toprule
Measure & Definition \\
\midrule
Intersection density & nodes with degree $\ge 3$ per km$^2$ \\
Street density & km of street per km$^2$ \\
Mean street length & mean edge length (m) \\
Streets per node & mean node degree (parallel edges counted) \\
Circuity & $\sum$ edge length $/ \sum$ straight-line length between edge endpoints \\
Dead-end, 3-way, 4-way shares & share of nodes with degree 1, 3 and 4 \\
Cyclomatic number & $e - v + p$ (independent loops; network redundancy) \\
Meshedness & $(e - v + p)/(2v - 5)$ (share of the maximum possible number of loops) \\
Gamma & $e / 3(v - 2)$ \\
Fragmentation $F$ & $1 - \operatorname{corr}(\ln L_i, G_i)$, see text \\
\bottomrule
\end{tabular}
\end{table}

\paragraph{Configurational fragmentation.} Space syntax diagnoses a ``patchwork'' city by the weak correlation
(low \emph{synergy}) between the local and global integration of its street segments \citep{medeiros2013urbis}.
Axial maps are not available for \nCities{} cities, so we compute a metric analogue on the OSM graph. For a node
$i$, local reach $L_i$ is the total street length whose two endpoints lie within 800~m of network distance from
$i$. Global closeness $G_i$ is the reciprocal of the mean network distance from $i$ to 500 randomly drawn urban
nodes. Synergy is the Pearson correlation between $\ln L_i$ and $G_i$ over 2,000 randomly drawn urban nodes,
and fragmentation is $F = 1 - \text{synergy}$. Routing uses the full buffered network. High $F$ means that
locally well-connected places are not necessarily well connected to the city as a whole. We stress that $F$
approximates the space-syntax measure with metric distances rather than replicating it.

\subsection{Analytical strategy}

For RQ1 we describe the distributions of $H^{E}$, $H^{N}$, $\Delta H$ and $\%\Delta H$, their Pearson and Spearman
correlations and the change in the ranking of the most segregated cities, and we repeat the comparison with
bandwidths of 1 and 3~km, an exponential kernel ($w = e^{-3d/b}$), the secondary indices and the
population-matched environments. For RQ2 we estimate, across cities,
\begin{equation}
  y_c = \alpha + \mathbf{x}_c'\boldsymbol\beta + \varepsilon_c ,
\end{equation}
with $y_c \in \{\Delta H, \%\Delta H\}$. $\mathbf{x}_c$ contains log intersection density, circuity, meshedness,
log cyclomatic number, dead-end share, log population density and $H^{E}$, plus the cyclomatic$\times$circuity
interaction of \citet{knaap2024segregated}. We then add $F$ and compare the fit. Normally distributed regressors
are standardised and skewed ones logged. Standard errors are heteroskedasticity-robust (HC3). With
$N = \nCities{}$ the models are descriptive associations, not causal estimates.
